% GENERATED FILE. Edit canonical lessons, then run npm run generate:books.
% canonical-source-hash: fnv1a64:360ce5e5cbff72e9
% canonical-lessons: IT-C195-pilota, IT-C195-commercialista, IT-C195-scienziato, IT-C195-ricercatore, IT-C195-turista

\chapter{Pilot, Accountant, Scientist, Researcher, Tourist}
\label{ch:it-pilot-accountant-scientist-researcher-tourist}
\hlchaptermodality{195}

\begin{chapteropening}
\textbf{By the end of this chapter:} \emph{I can say a pilot, an accountant, a scientist, a researcher and a tourist.}

Complete the last lesson of chapter 195: I can say a pilot, an accountant, a scientist, a researcher and a tourist.
\end{chapteropening}

\section[il pilota]{il pilota — a pilot}

\label{lesson:IT-C195-pilota}



\begin{quote}
Before the new one: say the Italian for an actress, then the Italian for a novel.
\end{quote}

\subsection*{You'll want to know: il pilota}
\textbf{il pilota} — \textquotedblleft{}a pilot\textquotedblright{}.

\begin{cousinweb}
\textbf{pilota} is usually traced to Byzantine Greek \emph{pedotes}, \textquotedblleft{}a steersman\textquotedblright{}, from \emph{pedon}, \textquotedblleft{}an oar blade\textquotedblright{}. EtymDB 2.1 has no row for its origin.
\end{cousinweb}

\begin{grammarlens}[title={What you've built}]
One more word for this chapter.
\end{grammarlens}

\subsection*{Your turn}
\begin{itemize}
\raggedright
  \item \emph{Say it:} \emph{il pilota}
  \item \emph{Say it:} \emph{il pilota}, once more
  \item \emph{Say it:} say \emph{il pilota}
  \item \emph{Recall:} say the Italian for an actress, then the Italian for a novel, then say \emph{il pilota} again
\end{itemize}

\begin{tcolorbox}[breakable,colback=teal!4,colframe=teal!35!black,title={Before you move on}]
What does il pilota mean? (\textquotedblleft{}A pilot\textquotedblright{}.) Say it once more.
\end{tcolorbox}
\section[il commercialista]{il commercialista — an accountant}

\label{lesson:IT-C195-commercialista}



\begin{quote}
Before the new one: say the Italian for a novel, then the Italian for a pilot.
\end{quote}

\subsection*{You'll want to know: il commercialista}
\textbf{il commercialista} — \textquotedblleft{}an accountant\textquotedblright{}.

\begin{cousinweb}
\textbf{commercialista} — built from \emph{commerciale}, \textquotedblleft{}of trade\textquotedblright{}, with \emph{-ista}. \emph{commercio}, \textquotedblleft{}trade\textquotedblright{} $\leftarrow$ Latin \emph{commercium}, from \emph{com-}, \textquotedblleft{}with\textquotedblright{}, and \emph{merx}, \textquotedblleft{}goods\textquotedblright{}. (EtymDB 2.1)
\end{cousinweb}

\begin{grammarlens}[title={What you've built}]
One more word for this chapter.
\end{grammarlens}

\subsection*{Your turn}
\begin{itemize}
\raggedright
  \item \emph{Say it:} \emph{il commercialista}
  \item \emph{Say it:} \emph{il commercialista}, once more
  \item \emph{Say it:} say \emph{il commercialista}
  \item \emph{Recall:} say the Italian for a novel, then the Italian for a pilot, then say \emph{il commercialista} again
\end{itemize}

\begin{tcolorbox}[breakable,colback=teal!4,colframe=teal!35!black,title={Before you move on}]
What does il commercialista mean? (\textquotedblleft{}An accountant\textquotedblright{}.) Say it once more.
\end{tcolorbox}
\section[lo scienziato]{lo scienziato — a scientist}

\label{lesson:IT-C195-scienziato}



\begin{quote}
Before the new one: say the Italian for a pilot, then the Italian for an accountant.
\end{quote}

\subsection*{You'll want to know: lo scienziato}
\textbf{lo scienziato} — \textquotedblleft{}a scientist\textquotedblright{}.

\begin{cousinweb}
\textbf{scienziato} — built from \emph{scienza}, \textquotedblleft{}science\textquotedblright{}, with \emph{-ato}. \emph{scienza} $\leftarrow$ Latin \emph{scientia}, \textquotedblleft{}knowledge\textquotedblright{}. (EtymDB 2.1)
\end{cousinweb}

\begin{grammarlens}[title={What you've built}]
One more word for this chapter.
\end{grammarlens}

\subsection*{Your turn}
\begin{itemize}
\raggedright
  \item \emph{Say it:} \emph{lo scienziato}
  \item \emph{Say it:} \emph{lo scienziato}, once more
  \item \emph{Say it:} say \emph{lo scienziato}
  \item \emph{Recall:} say the Italian for a pilot, then the Italian for an accountant, then say \emph{lo scienziato} again
\end{itemize}

\begin{tcolorbox}[breakable,colback=teal!4,colframe=teal!35!black,title={Before you move on}]
What does lo scienziato mean? (\textquotedblleft{}A scientist\textquotedblright{}.) Say it once more.
\end{tcolorbox}
\section[il ricercatore]{il ricercatore — a researcher}

\label{lesson:IT-C195-ricercatore}



\begin{quote}
Before the new one: say the Italian for an accountant, then the Italian for a scientist.
\end{quote}

\subsection*{You'll want to know: il ricercatore}
\textbf{il ricercatore} — \textquotedblleft{}a researcher\textquotedblright{}.

\begin{cousinweb}
\textbf{ricercatore} — built from \emph{ricercare}, \textquotedblleft{}to search again, to research\textquotedblright{}, with \emph{-tore}, \textquotedblleft{}one who\textquotedblright{}. \emph{ricercare} is \emph{ri-} plus \emph{cercare}, \textquotedblleft{}to look for\textquotedblright{} $\leftarrow$ Late Latin \emph{circare}. (EtymDB 2.1)
\end{cousinweb}

\begin{grammarlens}[title={What you've built}]
One more word for this chapter.
\end{grammarlens}

\subsection*{Your turn}
\begin{itemize}
\raggedright
  \item \emph{Say it:} \emph{il ricercatore}
  \item \emph{Say it:} \emph{il ricercatore}, once more
  \item \emph{Say it:} say \emph{il ricercatore}
  \item \emph{Recall:} say the Italian for an accountant, then the Italian for a scientist, then say \emph{il ricercatore} again
\end{itemize}

\begin{tcolorbox}[breakable,colback=teal!4,colframe=teal!35!black,title={Before you move on}]
What does il ricercatore mean? (\textquotedblleft{}A researcher\textquotedblright{}.) Say it once more.
\end{tcolorbox}
\section[il turista]{il turista — a tourist}

\label{lesson:IT-C195-turista}



\begin{quote}
Before the new one: say the Italian for a scientist, then the Italian for a researcher.
\end{quote}

\subsection*{You'll want to know: il turista}
\textbf{il turista} — \textquotedblleft{}a tourist\textquotedblright{}.

\begin{cousinweb}
\textbf{turista} $\leftarrow$ English \emph{tourist}, from \emph{tour} — borrowed. EtymDB 2.1 has no row for its origin.
\end{cousinweb}

\begin{grammarlens}[title={What you've built}]
One more word for this chapter.
\end{grammarlens}

\subsection*{Your turn}
\begin{itemize}
\raggedright
  \item \emph{Say it:} \emph{il turista}
  \item \emph{Say it:} \emph{il turista}, once more
  \item \emph{Say it:} say \emph{il turista}
  \item \emph{Recall:} say the Italian for a scientist, then the Italian for a researcher, then say \emph{il turista} again
\end{itemize}

\begin{tcolorbox}[breakable,colback=teal!4,colframe=teal!35!black,title={Before you move on}]
What does il turista mean? (\textquotedblleft{}A tourist\textquotedblright{}.) Say it once more.
\end{tcolorbox}
